\resizebox{\textwidth}{!}{%
\begin{tikzpicture}[
  scale=0.9,
  sub/.style={fill=brown!20},
  act/.style={fill=red!28},
  clad/.style={fill=green!20},
  mir/.style={fill=blue!20},
  phc/.style={fill=purple!20},
  elec/.style={fill=gray!65},
  outa/.style={->,>=Stealth,thick,orange!85!red},
  fba/.style={<->,thick,blue!70!black,dashed}
]
  \begin{scope}[shift={(0,0)}]
    \node[font=\bfseries] at (2.4,3.8) {(a) 边发射激光器};
    \draw[sub] (0,0) rectangle (4.8,0.75);
    \draw[clad] (0,0.75) rectangle (4.8,1.35);
    \draw[act] (0,1.35) rectangle (4.8,1.55);
    \draw[clad] (0,1.55) rectangle (4.8,2.15);
    \draw[thick] (0,0) -- (0,2.15);
    \draw[thick] (4.8,0) -- (4.8,2.15);
    \draw[fba] (0.5,2.45) -- (4.3,2.45) node[midway,above] {法布里-珀罗腔反馈};
    \draw[outa] (5.1,1.45) -- (6.2,1.45) node[midway,above] {边向出光};
    \node[below,font=\scriptsize] at (0,-0.2) {解理面};
    \node[below,font=\scriptsize] at (4.8,-0.2) {解理面};
    \node[font=\scriptsize] at (2.4,0.38) {衬底};
    \node[font=\scriptsize] at (2.4,1.05) {下包层};
    \node[font=\scriptsize] at (2.4,1.45) {有源区};
    \node[font=\scriptsize] at (2.4,1.85) {上包层};
  \end{scope}

  \begin{scope}[shift={(7.8,0)}]
    \node[font=\bfseries] at (2.4,3.8) {(b) 垂直腔面发射激光器};
    \draw[mir] (0,0) rectangle (4.8,1.1);
    \foreach \y in {0.12,0.30,...,1.0}{\draw[white,line width=1.5] (0,\y)--(4.8,\y);}
    \draw[act] (0,1.1) rectangle (4.8,1.45);
    \draw[mir] (0,1.45) rectangle (4.8,2.55);
    \foreach \y in {1.58,1.76,...,2.45}{\draw[white,line width=1.5] (0,\y)--(4.8,\y);}
    \draw[fill=gray!45] (1.8,1.1) rectangle (3.0,1.45);
    \draw[fba] (2.4,0.2) -- (2.4,2.35);
    \draw[outa] (2.4,2.7) -- (2.4,3.55) node[midway,right] {法向出光};
    \node[right,font=\scriptsize] at (4.95,0.55) {下反射镜};
    \node[right,font=\scriptsize] at (4.95,1.28) {有源区};
    \node[right,font=\scriptsize] at (4.95,2.0) {上反射镜};
    \node[font=\scriptsize] at (2.4,1.28) {限流孔};
    \draw[elec] (-0.25,2.3) rectangle (0,2.55);
    \draw[elec] (4.8,2.3) rectangle (5.05,2.55);
  \end{scope}

  \begin{scope}[shift={(0,-4.9)}]
    \node[font=\bfseries] at (2.4,3.8) {(c) 分布反馈激光器};
    \draw[sub] (0,0) rectangle (4.8,0.6);
    \draw[clad] (0,0.6) rectangle (4.8,1.0);
    \draw[phc] (0,1.0) rectangle (4.8,1.3);
    \foreach \x in {0.25,0.6,...,4.45}{\draw[fill=white,white] (\x,1.0) rectangle ++(0.16,0.3);}
    \draw[clad] (0,1.3) rectangle (4.8,1.85);
    \draw[act] (0,1.85) rectangle (4.8,2.05);
    \draw[fba] (0.5,2.35) -- (4.3,2.35) node[midway,above] {一维光栅反馈};
    \draw[out=270,in=90,->,thick,dotted] (2.4,1.3) to (2.4,0.72);
    \draw[outa] (5.1,1.45) -- (6.2,1.45) node[midway,above] {边向出光};
    \node[right,font=\scriptsize] at (4.95,0.3) {衬底};
    \node[right,font=\scriptsize] at (4.95,1.15) {一维光栅};
    \node[right,font=\scriptsize] at (4.95,1.95) {有源区};
  \end{scope}

  \begin{scope}[shift={(7.8,-4.9)}]
    \node[font=\bfseries] at (2.4,3.8) {(d) 光子晶体面发射激光器};
    \draw[sub] (0,0) rectangle (4.8,0.6);
    \draw[clad] (0,0.6) rectangle (4.8,1.05);
    \draw[act] (0,1.05) rectangle (4.8,1.35);
    \draw[phc] (0,1.35) rectangle (4.8,1.72);
    \foreach \x in {0.45,0.95,...,4.35}{
      \foreach \y in {1.42,1.54,1.66}{
        \draw[fill=white] (\x,\y) circle (0.1);
      }
    }
    \draw[clad] (0,1.72) rectangle (4.8,2.2);
    \draw[fba] (1.0,1.56) -- (3.8,1.56);
    \draw[fba] (2.4,1.40) -- (2.4,1.72);
    \node[font=\scriptsize,above] at (2.4,1.72) {二维面内反馈};
    \draw[outa] (2.4,1.8) -- (2.4,2.8) node[midway,right] {法向出光};
    \node[right,font=\scriptsize] at (4.95,0.3) {衬底};
    \node[right,font=\scriptsize] at (4.95,0.82) {下包层};
    \node[right,font=\scriptsize] at (4.95,1.2) {有源波导};
    \node[right,font=\scriptsize] at (4.95,1.54) {二维光子晶体层};
    \node[right,font=\scriptsize] at (4.95,1.95) {上包层/接触};
  \end{scope}

  \node[draw,rounded corners,fill=gray!10,align=left,anchor=west,text width=12.5cm]
    at (-0.3,-7.25) {\small\textbf{图注：}
    四种器件的核心差异在于反馈维度（一维或二维）与出光方向（边向或法向）。
    光子晶体面发射激光器的关键是“二维面内反馈 + 法向面发射”的统一实现。};
\end{tikzpicture}%
}
